\documentclass[11pt]{article}
\usepackage{icmlko}
\begin{document}

\title{검색 추이는 극단만 산다 --- 접전을 사는 자료가 무엇인지 처음으로 지목한다}
\author{월드모델 실험실}
\date{노트 537 · 2026년 8월 4일}
\maketitle

\begin{abstract}
노트 536 이 \textbf{채택 넷이 실무의 판($<2$ 배)을 거의 안 산다}는 것을 보였고,
\texttt{DOMDROP}(애니 검색 추이 가림)이 $5$ 배 이상에서만 버는 것에 대해
``검색 추이는 크게 흥한 것만 잡는 신호일 것''이라는 가설을 남겼다.
축 계열을 하나씩 끄고 구간별 손해를 쟀다.

\begin{center}
\begin{tabular}{lrrrr}
\hline
계열 & \textbf{$<2$ 배} & $2\!\sim\!5$ 배 & $5$ 배 이상 & 비 \\
\hline
\textbf{ALL5}(기획 속성 다섯) & $\mathbf{+0.00733}$ & $+0.02059$ & $+0.03019$ & $4.1$ \\
trend\_(검색 추이) & $+0.00317$ & $+0.00950$ & $+0.02005$ & $\mathbf{6.3}$ \\
tag\_ & $+0.00292$ & $+0.00938$ & $+0.00864$ & $3.0$ \\
\textbf{wiki\_} & $+0.00231$ & $+0.00392$ & $+0.00290$ & $\mathbf{1.3}$ \\
계절 & $+0.00142$ & $+0.00364$ & $+0.00396$ & $2.8$ \\
gen & $+0.00052$ & $+0.00122$ & $+0.00208$ & $4.0$ \\
grp & $+0.00052$ & $+0.00161$ & $+0.00106$ & $2.0$ \\
\hline
\end{tabular}
\end{center}

사전 등록 ①(\emph{trend\_ 의 $5$ 배$\uparrow$/$<2$ 배 $\ge 2$}) \textbf{통과}
($6.3$ 배). ②(\emph{$<2$ 배 $1$ 등이 trend\_ 가 아니다}) \textbf{통과} ---
\textbf{ALL5}, 곧 \emph{사람이 태깅한 기획 속성 다섯}이다.
\end{abstract}

\begin{figure}[t]
\centering
\includegraphics[width=\textwidth]{figs/f1.pdf}
\caption{\textbf{가} --- 계열별 구간 값어치. \textbf{나} --- 덮음당.
\textbf{다} --- 읽기.}
\end{figure}

\section{가설이 맞았다}

\textbf{검색 추이는 극단 전용이다} --- $5$ 배 이상에서 $+0.0201$ 인데 $<2$ 배
에서는 $+0.0032$ 로 $6.3$ 배 차이다. \texttt{DOMDROP} 이 왜 $5$ 배 이상에서만
벌고 실무의 판에서는 음수였는지가 설명된다(노트 536).

\textbf{반대쪽 끝은 위키다} --- 비가 $1.3$ 배로 거의 평평하다. 접전에서든
극단에서든 고르게 번다.

\section{덮음으로 나누면 그림이 또 바뀐다}

계열마다 채워진 비율이 크게 다르다. 덮음으로 나눈 $<2$ 배 값어치는

\begin{center}
\begin{tabular}{lrr}
\hline
계열 & 덮음당 $<2$ 배 & 덮음 \\
\hline
\textbf{wiki\_} & $\mathbf{+0.0330}$ & $7\%$ \\
trend\_ & $+0.0079$ & $40\%$ \\
ALL5 & $+0.0073$ & $100\%$ \\
tag\_ & $+0.0054$ & $54\%$ \\
계절 & $+0.0014$ & $100\%$ \\
gen $\cdot$ grp & $+0.0008$ & $66 \cdot 68\%$ \\
\hline
\end{tabular}
\end{center}

\textbf{위키가 $2$ 등의 $4.2$ 배다.} 덮음 $7\%$ 로 저만큼 번다.

\begin{quote}
\textbf{단서 --- 덮음당은 선형 외삽이다.} 위키를 $100\%$ 채우면 저만큼 번다는
보장이 없다. 지금 위키가 붙은 행이 \emph{원래 다른 종류}일 수도 있다(노트 489 가
같은 함정을 적었다 --- ``자료가 도왔거나 그 자료 붙는 행이 원래 다르거나'').
그래도 \textbf{어디를 먼저 시험할지}는 정해 준다.
\end{quote}

\section{수집 우선순위에 두 번째 축이 생겼다}

노트 508 은 \textbf{행당 판 기여}로 순위를 매겼다(모바일 $>$ 애니 $>$ 세계애니
$\cdots$). 여기는 \textbf{접전을 사는가}로 매긴다. \textbf{둘은 다른 물음이고
답도 다르다.}

\begin{center}
\begin{tabular}{ll}
\hline
물음 & 답 \\
\hline
어느 \emph{도메인}의 행을 늘리나 & 모바일 $>$ 애니 $>$ 세계애니(노트 508) \\
어느 \emph{자료 종류}가 접전을 사나 & \textbf{ALL5} $>$ trend\_ $>$ tag\_ $>$ wiki\_ \\
어느 종류가 \emph{덮음당} 접전을 사나 & \textbf{wiki\_} $>$ trend\_ $>$ ALL5 \\
\hline
\end{tabular}
\end{center}

\section{노트 530 과 맞물린다}

\begin{quote}
\textbf{접전을 제일 많이 사는 계열(ALL5)이 시장팝업에서는 셋만 채워지고 각각
두 값 $\cdot$ 두 값 $\cdot$ 일곱 값뿐이다}(노트 530). 시장팝업이 어느 셈법에서도
꼴찌인 것(노트 534)과 같은 자리다. \textbf{서로 다른 세 노트가 같은 곳을
가리킨다.}
\end{quote}

\section*{조항}

\begin{quote}
\textbf{자료의 값어치는 ``얼마나 버나''가 아니라 ``어느 판을 사나''로 적는다.}
같은 크기의 이득이라도 극단에서 번 것과 접전에서 번 것은 제품에서 값이 다르다.
\end{quote}

\section*{남은 것}

\begin{itemize}
\item \textbf{위키를 더 붙이면 정말 접전을 사나} --- 덮음 $7\% \to$ 더 채우는
실험. 지금은 선형 외삽이라 \emph{예보}일 뿐이다.
\item \textbf{ALL5 안에서 어느 축이 접전을 사나} --- 다섯을 하나씩 끄면 나온다.
시장팝업 태깅에서 무엇부터 채울지가 정해진다.
\item \textbf{접전 전용 자를 관문에 넣을지}(노트 $528 \cdot 532 \cdot 536$ 이
세 번 물었다) --- 넣으면 \texttt{DOMDROP} 을 재심사해야 한다.
\end{itemize}

\end{document}
